Three gaps motivate the broader agenda. \textbf{Systematic SAE analysis for pure vision models}: SAEs have rarely been studied on ViTs trained without language supervision; visual tokens encode image patches rather than words, and it remains an open empirical question whether sparse directions recover meaningful visual structure. \textbf{Spatial structure and class-token aggregation}: unlike language transformers, ViTs mix all patch tokens bidirectionally \cite{Raghu2021}, and tracing how patch evidence is integrated into the \texttt{[CLS]} token requires attention-level causal analyses that are still underdeveloped in the vision setting. \textbf{Cross-modal labelling}: when no internal text-image alignment exists, assigning interpretable labels to SAE features requires an external grounding mechanism.

In this project we directly address the first and third gaps with an SAE-based pipeline plus external CLIP grounding. The second gap is treated as important future work: our implementation does not perform attention-head patching or explicit \texttt{[CLS]}-aggregation tracing. Accordingly, the primary research question here is whether SAEs recover interpretable, monosemantic features from a purely visual ViT and whether those features show consistent causal effects under dataset-level ablation.
